\section{Results}
\label{sec:results}

% ROUGH OUTLINE -- one bullet = the expected finding (roughly one paragraph +
% one float in the final text). No numbers invented; all quantitative claims
% filled once results exist.

\subsection{Twin-experiment baseline}
\label{sec:results:twin}
\begin{itemize}
  \item Each solver against itself (\invcrime\ reference): forcing parameters
        recovered accurately with tight posteriors, fast monotone convergence,
        observation fit at the noise floor. The ``it works'' ceiling every
        cross-model result is read against. [Fig.: parameter convergence per
        solver.]
\end{itemize}

\subsection{Cross-model parameter estimation}
\label{sec:results:cross}
\begin{itemize}
  \item \textbf{Headline result.} Skill degrades relative to the \twinexp\
        once \truthmodel\ $\neq$ \assimmodel\ --- the signature of the
        \invcrime\ removed. [Fig.: twin vs.\ cross error per parameter;
        Table: \RMSE/\CRPS\ for all pairings.]
  \item Split identifiability: \inflowangle\ and \velmag\ remain recoverable
        with modest bias; \shearexp\ and \sgs\ become systematically biased
        (confident but wrong).
  \item Pairings are asymmetric (which solver is truth vs.\ assimilation
        matters); degradation expected to grow with \discrepancy\ ``distance''
        (LES$\to$LES vs.\ LES$\leftrightarrow$\lbmmodel).
\end{itemize}

\subsection{Observation-space vs.\ parameter-space skill}
\label{sec:results:obsvsparam}
\begin{itemize}
  \item \textbf{Key finding.} Cross-model runs can fit the sensors nearly at
        the noise floor while parameter error stays large --- good data
        mismatch masks wrong parameters, so obs-space residuals alone cannot
        flag failure. [Fig.: obs-space vs.\ parameter-space error scatter.]
\end{itemize}

\subsection{Role of the model-error knobs}
\label{sec:results:knobs}
\begin{itemize}
  \item Contrast three estimation targets (forcing only / $+\shearexp$ /
        $+\shearexp+\sgs$) on one cross-model pairing: freeing the knobs
        improves the observation fit by absorbing \discrepancy, but pushes the
        knobs to non-physical values and can degrade the forcing estimates.
        [Fig.: obs fit and parameter error across the three targets.]
\end{itemize}

\subsection{Sensitivity study}
\label{sec:results:sensitivity}
\begin{itemize}
  \item One factor at a time from the cross-model baseline; trim to the
        factors that matter once results are in. Expected: \Ne\ plateau;
        diminishing returns in \Na; optimum \Cd\ inflation balancing overfit
        vs.\ under-informing; localization helps at small \Ne; sensor
        placement beats count; prior spread trades bias vs.\ convergence.
        [Fig.: multi-panel skill vs.\ each retained factor.]
\end{itemize}

\subsection{Idealized vs.\ realistic geometry}
\label{sec:results:cases}
\begin{itemize}
  \item Does the twin$\to$cross degradation widen from the Xie \& Castro array
        to the Barcelona geometry, and which parameters are most affected?
        [Fig.: twin-vs-cross degradation, both cases side by side.]
\end{itemize}
